\subsection*{Actual pants insertions and coprime averaging}

\statement{B.54}{Proved}{An actual relative automorphism with nonuniform tripod actions}
There is a continuous peripheral automorphism of
$\widehat B/R_{\rm nil}(\widehat K)$, with multiplier one and
identity action on $A_5$, whose every full R2 lift lies outside
$\mathcal K$ and outside the Galois image. On the class-four,
exponent-seven quotients of the $220$ tripod groups of B.10,
its outer actions have degree-three coefficients
\[
       0\quad\hbox{on }160\hbox{ tripods},\qquad
       10q_3\ne0\quad\hbox{on }60\hbox{ tripods},             \tag{B.54.1}
\]
where $q_3$ is one specified nonzero arithmetic class over
$\mathbb F_7$. Consequently R2 does not admit a universal
strengthening requiring its full lifts to lie in the actual
profinite mapping class image.

The construction has the following more general form, used in
B.55. Let $U\trianglelefteq B$ be a torsion-free normal subgroup
of finite index, put $Q=B/U$, suppose $7\nmid |Q|$, and suppose
the surface of $U$ has at least three cusps. Choose two cusps
and a disc containing exactly those two punctures. There is
an actual pure peripheral automorphism
\[
 h\in\operatorname{Aut}(U^{(7)})
\]
which centralizes a finite complement $Q$ in
$\widehat B/O^7(\widehat U)$, is the identity modulo
$\gamma_4(U^{(7)})$, and whose degree-three derivation modulo
seven is the sum of all $Q$-translates of the specified pants
insertion. It extends to a peripheral automorphism of
$\widehat B/O^7(\widehat U)$ inducing the identity on $Q$.
Taking this automorphism at seven and the identity at every
other prime gives an R2 input at $U$. All full R2 lifts have
the same specified derivation on $U^{(7)}$.

\begin{proof}
We construct an actual arithmetic automorphism first, insert
it in an actual free pro-seven group, and only then use its
finite graded quotient. No assertion that an arbitrary Lie
derivation integrates is used.

\paragraph{A finite Kummer specification of the arithmetic seed.}
Fix $\zeta=\zeta_7$ and put
\[
 k_7=\mathbb Q(\zeta),\qquad
 \epsilon=\prod_{a=1}^6(1-\zeta^a)^{a^2},\qquad
 L=k_7(\epsilon^{1/7}).                                    \tag{B.54.2}
\]
The prime of $k_7$ above $29$ specified by $\zeta\mapsto16$
has residue field $\mathbb F_{29}$. Direct computation gives
\[
             \epsilon\mapsto21,\qquad21^4=7\ne1
                           \quad\hbox{in }\mathbb F_{29}.
\]
As $16$ has order seven and none of the factors vanishes,
$\epsilon$ is not a seventh power in $k_7$. Thus
$\operatorname{Gal}(L/k_7)=C_7$.
For $b\in\mathbb F_7^\times$, reindexing the factors in
(B.54.2) gives
\[
             b(\epsilon)\equiv\epsilon^{\,b^{-2}}
                              \pmod{(k_7^\times)^7}.
\]
Hence $L/\mathbb Q$ is Galois. Its quotient
$\operatorname{Gal}(k_7/\mathbb Q)$ acts on its normal $C_7$
by $b^3$: the root-of-unity action contributes $b$, and the
Kummer line contributes the inverse of $b^{-2}$.
For $b=3$ this action is inversion. Its commutator subgroup
is therefore exactly $C_7$, and its maximal abelian
subextension is $k_7$. Since $\mathbb Q(\mu_\infty)/\mathbb Q$
is abelian,
\[
                 L\cap\mathbb Q(\mu_\infty)=k_7.
\]
Choose, once and for all, an element
\[
 \sigma\in G_{\mathbb Q(\mu_\infty)},\qquad
 \sigma(\epsilon^{1/7})=\zeta\,\epsilon^{1/7}.               \tag{B.54.3}
\]
The intersection equality proves the existence of this
element; (B.54.3) is a literal finite extension condition,
rather than an unspecified nonzero regulator.

Write $\kappa_3$ for the Soul\'e character normalized by the
Kummer products in (B.54.2). Formula (1.2) of
\cite[pp.~55--56]{CENakamuraSakugawaWojtkowiak2017} gives
$\kappa_3(\sigma)\equiv1\pmod7$.
The depth-one coefficient formula in
\cite[\S5, pp.~718--719]{CEFurusho2003} identifies the
degree-three coefficient of its Ihara element with
\[
             \frac{\kappa_3(\sigma)}{2(7^2-1)}
\]
up to the harmless overall sign from the convention for the
Ihara action. In particular this coefficient is a
seven-adic unit. These are the only arithmetic coefficient
formulas used here; neither the generation conjecture for
the Galois Lie algebra nor an equality with the full
Grothendieck--Teichm\"uller Lie algebra is needed.

Let $P_2$ be the free pro-seven group on $x,y$ and
$t=xy$. The multiplier of $\sigma$ is one. Its standard
based representative $g_0$ fixes $x$ and sends
$y$ to $f^{-1}yf$. The degree-two coefficient of $f$ is
$(\lambda^2-1)/24=0$, so $f\in\gamma_3(P_2)$ and
$g_0$ is the identity modulo $\gamma_4(P_2)$.
The two-cycle relation makes the two degree-three
coefficients equal. If $X,Y$ are the degree-one variables
and $C=[X,Y]$, its degree-three term consequently is
\[
                  f_3=a\bigl([X,C]+[Y,C]\bigr),
                          \qquad a\in\mathbb Z_7^\times.  \tag{B.54.4}
\]

We need a representative fixing $t$ exactly. Write
$g_0(t)=ktk^{-1}$. The congruence $[k,t]\in\gamma_4(P_2)$
first says that the degree-one part of $k$ is a multiple
of $X+Y$. Multiplying $k$ on the right by the corresponding
power of $t$ removes it without changing $ktk^{-1}$.
The degree-two part then centralizes $X+Y$ in the free
Lie algebra, and is zero: a multiple of $[X,Y]$ cannot
commute with $X+Y$. Thus $k$ may be chosen in
$\gamma_3(P_2)$. The automorphism
$g=\operatorname{Inn}(k^{-1})g_0$ fixes $t$ and remains
the identity modulo $\gamma_4(P_2)$. It is still an actual
continuous automorphism of $P_2$ preserving the two
puncture classes.

\paragraph{The exact first nonzero derivation.}
All the following Lie expressions are in degree at most
four over $\mathbb F_7$. Put
\[
 A=[X,[X,Y]],\quad B=[Y,[X,Y]],\quad
 U_4=[X,A],\quad V_4=[Y,A]=[X,B],\quad W_4=[Y,B].
\]
The first derivation of the boundary-fixed representative
is
\[
               d_{12}(X)=-aV_4,\qquad
               d_{12}(Y)=aV_4.                            \tag{B.54.5}
\]
Indeed $g_0$ has derivation $(0,a(V_4+W_4))$.
The degree-three part of $k$ is $-aB$, and the indicated
inner adjustment gives (B.54.5).
Let $q_3$ be the outer class of (B.54.5).
It is nonzero. For an inner derivation in this degree the
conjugating element is $\alpha A+\beta B$ and its values
on $X,Y$ are
\[
           (-\alpha U_4-\beta V_4,\,
                            -\alpha V_4-\beta W_4).
\]
Comparison with $a(-V_4,V_4)$ would force
$\alpha=0,\ \beta=a$, followed by
$-aW_4=aV_4$, which is impossible in the free Lie algebra.
Inner elements of smaller degree cannot intervene: their
first nonzero homogeneous term would have to centralize
both $X$ and $Y$.

For later changes of the named pair let $Z=-X-Y$.
The insertions for the pairs $\{X,Z\}$ and $\{Y,Z\}$
have derivations
\[
 d_{13}=a(-U_4-V_4,0),\qquad
 d_{23}=a(0,V_4+W_4).
\]
Their differences from $d_{12}$ are respectively
$a\,\operatorname{ad}(A)$ and
$-a\,\operatorname{ad}(B)$. Thus all three pairs determine
the same nonzero outer class $q_3$. Interchanging the two
members of a pair also leaves the insertion derivation
unchanged. This verifies directly the needed invariance
under all three-cusp permutations.

These computations are genuine finite-group tests. For a
free pro-seven group $P$ of finite rank use
\[
       N_4(P)=P/\overline{P^7\gamma_5(P)}.                 \tag{B.54.6}
\]
It is the free nilpotent group of class at most four and
exponent seven on the given generators. Its logarithm is
the free Lie algebra through degree four over
$\mathbb F_7$; the truncated exponential, logarithm and
Campbell--Hausdorff formulas have denominators prime to
seven. This is the class-less-than-$p$ correspondence of
\cite[II, \S4, Theorems 4.2--4.3 and Lemma 4.5]{CELazard1954}.
One can also obtain this identification directly by
applying those polynomials in the truncated tensor algebra:
successive commutators span each homogeneous free Lie
piece, and the universal property gives (B.54.6).
In particular the displayed non-inner derivation is
detected in an actual finite quotient, not only after
tensoring with $\mathbb Q_7$.

\paragraph{Insertion and a fixed lift in the automorphism group.}
Put $P=U^{(7)}$. A standard surface presentation may be
chosen with the two specified puncture meridians $x,y$,
the remaining puncture meridians and the handle generators.
The disc boundary is $xy$. Apply $g$ to $x,y$ and the
identity to the fundamental group of the complementary
surface. Since $g(xy)=xy$, the amalgamation relation is
respected. Equivalently use a free basis omitting a third
puncture meridian: the product relation defining that
meridian is unchanged. The inverse is obtained from
$g^{-1}$. This gives an actual automorphism
$a_{ij}\in\operatorname{Aut}(P)$, pure peripheral and
trivial modulo $\gamma_4(P)$, in arbitrary genus.

Set $E=\widehat B/O^7(\widehat U)$. Schur--Zassenhaus
provides a complement $Q\leq E$ to its normal pro-seven
subgroup $P$. Conjugation by this complement is an actual
finite action on $P$. Its outer action is the usual
orientation-preserving deck action, so it permutes the
puncture classes with exponent one.
Let
\[
 H=\overline{\langle\,q a_{ij}q^{-1}:q\in Q\,\rangle}
                          \leq\operatorname{Aut}(P).      \tag{B.54.7}
\]
Here $q a_{ij}q^{-1}$ denotes conjugation by the
automorphism of $P$ supplied by $q$.
This is a $Q$-stable pro-seven group. To see the last
point, the kernel of
$\operatorname{Aut}(P)\to\operatorname{Aut}(P/\Phi(P))$
is pro-seven: its images on all finite characteristic
seven-group quotients are seven-groups by the finite
Burnside basis theorem. Every generator of (B.54.7) lies
in this kernel. Moreover $H$ consists of pure peripheral
automorphisms trivial modulo $\gamma_4(P)$; both conditions
are closed and are preserved by conjugation by $Q$.

Let $V$ be the image of $H$ in
$\operatorname{Aut}(N_4(P))$. This is a finite elementary
abelian seven-group. An automorphism in the image changes
each generator by a central element of degree four, so
composition adds its degree-three derivation; it fixes
the degree-four subgroup pointwise. Thus $V$ is exactly
a finite quotient on which all the required degree-three
data can be read. It has a $Q$-action. If $v_{ij}$ denotes
the image of $a_{ij}$, then
\[
                     v=\sum_{q\in Q}qv_{ij}\in V^Q.        \tag{B.54.8}
\]

For completeness the coprime fixed-point lifting used
here is an exact group argument. Write $W=\ker(H\to V)$,
choose $b\in H$ mapping to $v$, and work in $H\rtimes Q$.
Because $v$ is fixed, both $Q$ and $bQb^{-1}$ are
complements to $W$ in $W\rtimes Q$. Schur--Zassenhaus
conjugates these complements by an element $w\in W$.
Then $h=w^{-1}b$ normalizes $Q$. Its conjugation on $Q$
is the identity, since $h\in H$ and projection to $Q$
is unchanged. Consequently
\[
                  h\in C_H(Q),\qquad h\mapsto v.          \tag{B.54.9}
\]
The profinite form follows from the finite theorem by
compactness of the conjugator sets; its normal subgroup
is pro-seven and its quotient has order prime to seven.
This proves existence at the full pro-seven level,
without any lift of an arbitrary formal derivation.
In particular it preserves the entire trace (B.54.8),
and not just its images on a chosen collection of
tripods. On boundary homology classes the $Q$-action is
permutation, even when the surface has positive genus;
hence (B.54.5) computes every boundary value of the trace.

\paragraph{Extension, the cusp root, and R2.}
Since $h$ centralizes $Q$, extend it to
$\alpha=h\rtimes\operatorname{id}_Q$ on $E=P\rtimes Q$.
It sends $s,r$ to conjugates by elements of $P$. Indeed
each of these elements has order prime to seven, and
the complements in the preimage of its cyclic image
in $Q$ are $P$-conjugate. The chosen complement is
fixed by $\alpha$.

Let $m$ be the order of $\bar\tau$ in $Q$ and put
$c=\tau^m\in P$. This is a primitive puncture meridian;
$m$ is prime to seven. Since there are at least three
cusps, their distinct classes in $H_1(P,\mathbb Z_7)$
have exactly the stabilizers of the corresponding
cusps. Therefore any element of $C_E(c)$ projects to
the cusp stabilizer $\langle\bar\tau\rangle$.
The centralizer of the primitive element $c$ in the
free pro-seven group $P$ is $\langle c\rangle$.
After multiplication by a suitable power of $\tau$,
an element of $C_E(c)$ is in this latter centralizer.
We obtain the exact equality
\[
                C_E(c)=\overline{\langle\tau\rangle}
                               \simeq\mathbb Z_7\times C_m.
                                                               \tag{B.54.10}
\]
Here the centralizer assertion inside $P$ also follows
by writing $P=\langle c\rangle\amalg_7 P'$ and using
the free pro-seven product tree.
Because $h$ is pure peripheral, write $h(c)=ucu^{-1}$
with $u\in P$. The element
$b=u^{-1}\alpha(\tau)u$ satisfies $b^m=c$ and has
the same image as $\tau$ in $Q$. In (B.54.10) an
$m$th root with that prescribed finite image is unique,
since multiplication by $m$ is invertible on
$\mathbb Z_7$. Thus $b=\tau$.
All three required conjugators are in $P$.

Take $\alpha$ at seven and the identity at every other
prime in the fibre product
\[
 \widehat B/R_{\rm nil}(\widehat U)
       =\prod_{p,Q}\widehat B/O^p(\widehat U).
                                                               \tag{B.54.11}
\]
Their quotient actions agree, and their conjugators
also agree over $Q$: the seven-component is in $P$,
and all other components are the identity. This is
an actual continuous automorphism of (B.54.11) with
the three prescribed conjugacy conditions. R2 supplies
full lifts with these conjugators lifted. Their outer
classes are in $\mathcal H$ by B.6, and their action
on $U^{(7)}$ has exactly (B.54.8). This proves the
general construction.

\paragraph{The icosahedral trace and its consequences.}
Now take $U=K$ and $Q=A_5$. Its twelve cusps are the
cosets of $C_5$; the six antipodal pairs are the
blocks $D_5/C_5$ coming from the normalizer of $C_5$.
Insert the seed in one antipodal pair. In the
$60$-term sum (B.54.8), each of the six antipodal
pairs occurs exactly $10$ times.
Killing all but three cusp classes kills a pair
insertion unless both members of its pair survive.
If both survive, its outer first derivation is $q_3$,
independently of the names of the surviving cusps,
by the explicit inner-derivation calculation above.
A three-subset contains at most one antipodal pair.
There are $6\cdot10=60$ three-subsets containing
one, and $\binom{12}{3}-60=160$ containing none.
This proves (B.54.1), with $10q_3\ne0$ modulo seven.

Every element of $\mathcal K$ has trivial action on
each of these tripod groups by B.10. Thus none of the
full R2 lifts just constructed belongs to $\mathcal K$.
If such a lift were arithmetic, its multiplier one
would place it over $\mathbb Q(\mu_\infty)$, and the
geometrically marked tripod maps of B.10 would all
have its same arithmetic action. Their degree-three
classes would in particular be equal. The unequal
values in (B.54.1) contradict that conclusion.
Finally a lift in the actual profinite mapping class
image would be in $\mathcal K$ by its definition,
giving the stated counterexample to the universal
braid-constrained lifting assertion.
\end{proof}

The executable certificate
\texttt{work\_CE/code/check\_pants\_average.py}
checks the finite Kummer reduction, all three Lie
formulas, non-innerness, and all $60$ permutations and
$220$ triples used above. The existence of the full
automorphism is proved by (B.54.7)--(B.54.11), not
inferred from those finite checks. This construction
does not produce an element of $\mathcal K$ or impose
the pentagon on any of its R2 lifts.

